\documentclass[english, a4paper, 10pt]{ieej-tec}
%\usepackage[dvips]{graphicx}
\usepackage[dvipdfmx]{graphicx}
\usepackage[cmex10]{amsmath}

\setlength{\topmargin}{-20mm}
\setlength{\oddsidemargin}{-32mm}
\setlength{\evensidemargin}{-32mm}
%\pagestyle{empty}

\title{Paper Title}
\authorlist{%
 \authorentry*{Taro DENKI}{DU}
 \authorentry{Jiro DENKI}{DU}
 \authorentry{Hanako DENSHI}{DEU}
 \authorentry{Ichiro DENSHI}{DEU}
}
\affiliate[DU]{Denki University}
\affiliate[DEU]{Denshi University}

%¥breakaffiliate{1}
%¥breakJauthorline{3}
%¥breakEauthorline{3}

\begin{document}
\begin{abstract}
Most word limits are around 200 words.
\end{abstract}

\begin{keyword}
AVIC, ECT, Integrated circuit, Analog design, ..., (5--6 words).
\end{keyword}
\maketitle

\section{Introduction}
This template, using ieej-tec.cls for the PC, provides authors with most of the formatting specifications needed for preparing electronic versions of their papers. All standard paper components have been specified for three reasons: (1) ease of use when formatting individual papers, (2) automatic compliance to electronic requirements that facilitate the concurrent or later production of electronic products, and (3) conformity of style throughout a conference proceedings. Margins, column widths, line spacing, and type styles are built-in.


\section{Section 2}

\begin{figure}[b]
\begin{center}
\includegraphics[scale=0.7]{avic_logo.eps}
\caption{AVIC logo.}
\end{center}
\end{figure}

\begin{table}[b]
  \centering
  \caption{Device parameters.} \label{table1}
  \begin{tabular}{lcr}  
  \hline
    M1 [$\mu$m]  & W/L= 10/ 0.18\\
    M2 [$\mu$m]  & W/L= 60/ 0.18\\
    C				    & 0.25 pF\\ \hline
  \end{tabular}
\end{table}

\subsection{Subsection}

\section{Conclusion}

%===================================================================================
%References Here
%=================================================================================== 
\begin{thebibliography}{9}
\bibitem{CTIA}
S.M.\ Park and H.-J.\ Yoo, ``1.25-Gb/s regulated cascode CMOS transimpedance amplifier for gigabit ethernet applications,''
{\em IEEE J. Solid-state Circuit} vol.\ 39, no.\ 1, pp.\ 112--121, Jan.\ 2004.
\end{thebibliography}
\end{document}
